section{Baseline Analysis and Dataset Structure}

We begin by training a standard linear text classifier based on TF--IDF representations (word- and character-level n-grams), source indicators, and simple numeric features.
Hyperparameter tuning reveals a narrow performance plateau, with optimal configurations consistently converging to:
\[
\text{word\_ng} = 2,\quad \text{char\_ng} \in [5,7],\quad C \in [0.6,1.0],\quad \text{max\_df} \approx 0.85{-}0.9.
\]

Despite extensive exploration of the hyperparameter space, the macro-F1 score saturates around $0.719$, suggesting that the linear model has largely exhausted the information available under a single-stage formulation.
\subsection{Low-Entropy Lexical and Editorial Signals}

Inspection of feature weights reveals that the most influential predictors are not semantic tokens, but editorial and structural markers, including publisher identifiers, RSS-related tokens, and HTML or URL fragments.
Many of these features exhibit strong class-specific associations.

To quantify this behavior, we analyze token-level statistics and identify a subset of tokens with:
\begin{itemize}
	\item empirical class purity equal to 1,
	\item zero empirical entropy,
	\item non-negligible support.
\end{itemize}

Formally, let $t$ denote a token and $Y$ the class label.
The empirical entropy of $t$ is defined as:
\[
H(Y \mid t) = - \sum_{c \in \mathcal{Y}} p(c \mid t)\log p(c \mid t).
\]
For a non-trivial subset of tokens, we observe $H(Y \mid t) = 0$, indicating that their presence alone is sufficient to determine the class on the development set.
\subsection{Motivation}

The dataset exhibits a clear structural heterogeneity.
A significant fraction of documents contains near-deterministic editorial or lexical cues, while the remaining instances require genuinely probabilistic reasoning.

A single-stage classifier implicitly averages these two regimes, treating deterministic and ambiguous examples as if they originated from the same data-generating process.
This motivates an explicit decomposition of the problem into two stages, each aligned with a distinct inductive bias.
\section{Two-Stage Classification Model}

Let $\mathcal{D} = \{(x_i, y_i)\}_{i=1}^N$ be the dataset, where $x_i$ is a document and $y_i \in \mathcal{Y}$ its label.
Let $\mathcal{T}$ denote the vocabulary of extracted tokens.

\subsection{Stage 1: High-Purity Rule-Based Assignment}

For each token $t \in \mathcal{T}$, we compute:
\[
\text{purity}(t) = \max_{c \in \mathcal{Y}} p(c \mid t), \quad
\text{support}(t) = |\{x_i : t \in x_i\}|.
\]

We define a rule set:
\[
\mathcal{R} = \{t \in \mathcal{T} : \text{purity}(t) \ge \tau_p \;\wedge\; \text{support}(t) \ge \tau_s\},
\]
where $\tau_p$ and $\tau_s$ are fixed thresholds.

For a document $x$, if there exists $t \in \mathcal{R}$ such that $t \in x$, the document is assigned the label:
\[
\hat{y}^{(1)}(x) = \arg\max_{c \in \mathcal{Y}} p(c \mid t).
\]
\subsection{Stage 2: Probabilistic Classification}

Documents not covered by Stage 1 are passed to a standard probabilistic classifier.
Let $f_\theta(x)$ denote a linear model operating on TF--IDF features and auxiliary signals.
For uncovered documents, the final prediction is:
\[
\hat{y}^{(2)}(x) = \arg\max_{c \in \mathcal{Y}} f_\theta(x)_c.
\]
\subsection{Unified Decision Rule}

The overall classifier $\hat{y}(x)$ is defined as:
\[
\hat{y}(x) =
\begin{cases}
\hat{y}^{(1)}(x), & \text{if } \exists t \in \mathcal{R} \text{ s.t. } t \in x, \\
\hat{y}^{(2)}(x), & \text{otherwise}.
\end{cases}
\]
\subsection{Coverage--Performance Trade-off}

By varying $(\tau_p, \tau_s)$, we control the fraction of documents resolved in Stage 1.
Empirically, we observe that for $\tau_p \ge 0.95$ and $\tau_s \in [20,40]$, Stage 1 covers approximately $16\%$--$20\%$ of the dataset.

This early resolution improves the overall macro-F1 score while reducing the complexity of the decision boundary faced by Stage 2.
The performance gain saturates beyond this range, indicating that the most informative deterministic patterns have already been captured.
\subsection{Interpretability and Generality}

The proposed two-stage formulation does not rely on dataset-specific heuristics.
It is applicable to any classification task where:
\begin{itemize}
	\item a subset of features exhibits near-deterministic label associations,
	\item the remaining data requires probabilistic modeling,
	\item shortcut signals and genuine ambiguity coexist.
\end{itemize}

By explicitly separating these regimes, the model improves interpretability, robustness, and sample efficiency without increasing architectural complexity.
\section{Baseline Analysis and Model Motivation}

We first consider a standard linear text classification approach based on TF--IDF representations, source indicators, and simple numeric features.
Despite extensive exploration of the hyperparameter space, performance consistently converges to a narrow plateau, with optimal configurations characterized by:
\[
\text{word\_ng}=2,\quad \text{char\_ng}\in[5,7],\quad C\in[0.6,1.0],\quad \text{max\_df}\approx0.85{-}0.9.
\]
Across these configurations, the macro-averaged F1 score saturates around $0.719$, indicating that the linear model has largely exhausted the information available under a single-stage formulation.

Inspection of learned feature weights reveals that the most influential predictors are predominantly editorial and structural rather than semantic. These include publisher identifiers, RSS-related markers, and HTML or URL fragments. Such features exhibit strong class-specific associations and dominate the linear decision function.

To characterize this behavior, we analyze token-level statistics and identify a subset of tokens whose empirical class distribution exhibits zero entropy. Let $t$ denote a token and $Y$ the class label. The empirical conditional entropy is defined as:
\[
H(Y \mid t) = - \sum_{c \in \mathcal{Y}} p(c \mid t)\log p(c \mid t).
\]
For a non-trivial subset of tokens, we observe $H(Y \mid t)=0$, meaning that their presence alone determines the class on the development set with certainty. This evidences a structural heterogeneity in the dataset: a fraction of documents is governed by near-deterministic editorial cues, while the remainder requires genuinely probabilistic reasoning.

A single-stage classifier implicitly averages these two regimes, motivating an explicit architectural decomposition aligned with their distinct inductive biases.
\section{Two-Stage Classification Model}

Let $\mathcal{D}=\{(x_i,y_i)\}_{i=1}^N$ be the dataset, where $x_i$ is a document and $y_i\in\mathcal{Y}$ its label, and let $\mathcal{T}$ denote the vocabulary of extracted tokens.

\subsection{Stage 1: High-Purity Rule-Based Assignment}

For each token $t\in\mathcal{T}$, we compute:
\[
\text{purity}(t)=\max_{c\in\mathcal{Y}} p(c\mid t), \quad
\text{support}(t)=|\{x_i : t\in x_i\}|.
\]
We define a rule set:
\[
\mathcal{R}=\{t\in\mathcal{T} : \text{purity}(t)\ge\tau_p \;\wedge\; \text{support}(t)\ge\tau_s\}.
\]
If a document $x$ contains any token $t\in\mathcal{R}$, it is assigned the label:
\[
\hat{y}^{(1)}(x)=\arg\max_{c\in\mathcal{Y}} p(c\mid t).
\]

\subsection{Stage 2: Probabilistic Classification}

Documents not covered by Stage~1 are passed to a standard linear classifier $f_\theta(x)$ operating on TF--IDF features and auxiliary signals. For such documents, the prediction is:
\[
\hat{y}^{(2)}(x)=\arg\max_{c\in\mathcal{Y}} f_\theta(x)_c.
\]

\subsection{Unified Decision Rule}

The final prediction is defined as:
\[
\hat{y}(x)=
\begin{cases}
\hat{y}^{(1)}(x), & \text{if } \exists\, t\in\mathcal{R}\text{ such that } t\in x,\\
\hat{y}^{(2)}(x), & \text{otherwise}.
\end{cases}
\]
\subsection{Coverage and Robustness}

By varying $(\tau_p,\tau_s)$, the fraction of documents resolved in Stage~1 can be controlled.
Empirically, for $\tau_p\ge0.95$ and $\tau_s\in[20,40]$, Stage~1 resolves approximately $16\%$--$20\%$ of the dataset.
These instances correspond to documents containing near-deterministic editorial or structural markers.

The two-stage formulation is robust by construction.
In the worst-case scenario, where no document satisfies the rule constraints (e.g., due to domain shift or unseen tokens), the model degenerates gracefully to the Stage~2 classifier, recovering the single-stage linear baseline without performance loss.
Conversely, when deterministic patterns are present, they are exploited without contaminating the probabilistic decision boundary.

This separation improves interpretability and sample efficiency while avoiding overconfident assignments on ambiguous instances.
